% Author: Izaak Neutelings (October 2021)
% Inspiration
%   "Very special relativity - An illustrated guide", Sander Bais (2007)
%   http://people.uncw.edu/hermanr/GR/Minkowski/Minkowski.pdf
\documentclass[border=3pt,tikz]{standalone}
\usepackage{amsmath} % for \text
\usepackage{etoolbox} % ifthen
\usepackage[outline]{contour} % glow around text
\usetikzlibrary{calc} % for adding up coordinates
\usetikzlibrary{decorations.markings,decorations.pathmorphing}
\usetikzlibrary{angles,quotes} % for pic (angle labels)
\usetikzlibrary{arrows.meta} % for arrow size
\usepackage{xfp} % higher precision (16 digits?)
\contourlength{1.1pt}

\tikzset{>=latex} % for LaTeX arrow head
\colorlet{myred}{red!85!black}
\colorlet{mydarkred}{red!55!black}
\colorlet{mylightred}{red!85!black!12}
\colorlet{myfieldred}{mydarkred!5} % for S' background
\colorlet{myredhighlight}{myred!20} % highlights simultaneity in ladder paradox
\colorlet{myblue}{blue!80!black}
\colorlet{mydarkblue}{blue!50!black}
\colorlet{mylightblue}{blue!50!black!30}
\colorlet{mylightblue2}{myblue!10}
\colorlet{mygreen}{green!80!black}
\colorlet{mypurple}{blue!40!red!80!black}
\colorlet{mydarkgreen}{green!50!black}
\colorlet{mydarkpurple}{blue!40!red!50!black}
\colorlet{myorange}{orange!40!yellow!95!black}
\colorlet{mydarkorange}{orange!40!yellow!85!black}
\colorlet{mybrown}{brown!20!orange!90!black}
\colorlet{mydarkbrown}{brown!20!orange!55!black}
\colorlet{mypurplehighlight}{mydarkpurple!20} % highlights simultaneity in ladder paradox
\tikzstyle{world line}=[myblue!40,line width=0.3]
\tikzstyle{world line t}=[mypurple!50!myblue!40,line width=0.3]
\tikzstyle{world line'}=[mydarkred!40,line width=0.3]
\tikzstyle{mysmallarr}=[-{Latex[length=3,width=2]},thin]
\tikzstyle{mydashed}=[dash pattern=on 3 off 3]
\tikzstyle{rod}=[mydarkbrown,draw=mydarkbrown,double=mybrown,double distance=2pt,
                 line width=0.2,line cap=round,shorten >=1pt,shorten <=1pt]
%\tikzstyle{rod'}=[rod,draw=mydarkbrown!80!red!85,double=mybrown!80!red!85]
\tikzstyle{vector}=[->,line width=1,line cap=round]
\tikzstyle{vector'}=[vector,shorten >=1.2]
\tikzstyle{particle}=[mygreen,line width=0.9]
\tikzstyle{photon}=[-{Latex[length=5,width=4]},myorange,line width=0.8,decorate,
                    decoration={snake,amplitude=1.0,segment length=5,post length=5}]

\def\tick#1#2{\draw[thick] (#1) ++ (#2:0.06) --++ (#2-180:0.12)}
\def\tickp#1#2{\draw[thick,mydarkred] (#1) ++ (#2:0.06) --++ (#2-180:0.12)}
\def\Nsamples{100} % number samples in plot

\begin{document}

% COMMON AXES
\pgfdeclarelayer{back} % to draw on background
\pgfsetlayers{back,main} % set order
\def\xmin{2}
\def\xmax{2}
\def\Nlines{6} % number of world lines (at constant x/t)
\def\DNxp{0}   % difference in number of world lines of x' axis
\def\DNyp{0}   % difference in number of world lines of ct' axis
\def\DNy{0}    % difference in number of world lines of ct axis
\def\ang{20}   % angle between x and x' axes
\def\xplabelang{180} % anchor angle of x' axis label
%\pgfmathsetmacro\ang{atan(0.44)} % angle between x and x' axes
\def\axes{
  \pgfmathsetmacro\d{\xmax/(\Nlines+0.4)} % grid size
  \pgfmathsetmacro\D{\d/cos(\ang)/sqrt(1-tan(\ang)^2)} % grid size, boosted
  \pgfmathsetmacro\ymax{\xmax+\DNy*\d} % maximum of y = ct axis
  \pgfmathsetmacro\xmaxp{(\xmax/\d+\DNxp)*\D} % maximum of x' axis
  \pgfmathsetmacro\ymaxp{(\xmax/\d+\DNyp)*\D} % maximum of y' = ct' axis
  \pgfmathsetmacro\Nylines{\Nlines+\DNy} % number of world lines at constant ct'
  \pgfmathsetmacro\Nxplines{\Nlines+\DNxp} % number of world lines at constant x'
  \pgfmathsetmacro\Nyplines{\Nlines+\DNyp} % number of world lines at constant ct'
  \coordinate (O) at (0,0);
  \coordinate (X) at (\xmax+0.15,0);
  \coordinate (T) at (0,\ymax+0.15);
  \coordinate (X') at (\ang:\xmaxp+0.2);
  \coordinate (T') at (90-\ang:\ymaxp+0.2);
  
  % FILL
  % \begin{pgfonlayer}{back} % draw on back
  %   \fill[myfieldred]
  %     (\ang:-\xmin) -- (\ang:\xmaxp) --++ (90-\ang:\ymaxp) --++ (\ang:-\xmaxp)
  %     -- (90-\ang:-\xmin) -- cycle;
  % \end{pgfonlayer}
  
  % WORLD LINE GRID
  \message{  Making world lines...^^J}
  \foreach \i [evaluate={\x=\i*\d;}] in {-\Nlines,...,\Nlines}{
    %\message{  Running i/N=\i/\Nlines, x=\x...^^J}
    \draw[world line]   (\x,-\xmin) -- (\x,\ymax);
  }
  \foreach \i [evaluate={\t=\i*\d;}] in {-\Nlines,...,\Nylines}{
    %\message{  Running i/N=\i/\Nlines, t=\t...^^J}
    \draw[world line t] (-\xmin,\t) -- (\xmax,\t);
  }
  
  % BOOSTED WORLD LINE GRID
  % \message{  Making world lines for boosted frame...^^J}
  % \foreach \i [evaluate={\x=\i*\D;}] in {1,...,\Nxplines}{
  %   %\message{  Running i/N=\i/\Nlines, x=\x...^^J}
  %   \draw[world line'] (\ang:\x) --++ (90-\ang:\ymaxp);
  % }
  % \foreach \i [evaluate={\t=\i*\D;}] in {1,...,\Nyplines}{
  %   %\message{  Running i/N=\i/\Nlines, t=\t...^^J}
  %   \draw[world line'] (90-\ang:\t) --++ (\ang:\xmaxp);
  % }
  
  % AXES
  \draw[->,thick] (0,-\xmin) -- (T) node[above] {$ct$};
  \draw[->,thick] (-\xmin,0) -- (X) node[right] {$x$};
  % \draw[->,thick,mydarkred] (90-\ang:-\xmin) -- (T')
    % node[right=5,above=-1] {$ct'$};
  % \draw[->,thick,mydarkred] (\ang:-\xmin) -- (X')
    % node[anchor=\xplabelang,inner sep=2] {$x'$};
}



% SPACETIME DIAGRAM - LENGTH CONTRACTION of rod at rest in S
% Inspiration: http://people.uncw.edu/hermanr/GR/Minkowski/Minkowski.pdf
\def\ang{23} % angle between x and x' axes
\begin{tikzpicture}[scale=1.8]
  \message{Length contraction (rod at rest in S)^^J}
  
  % AXES
  \def\Nlines{7} % number of world lines (at constant x/t)
  \axes
  
  % SETTINGS
  \pgfmathsetmacro\xA{0} % triangle left corner x coordinate in S
  \pgfmathsetmacro\yA{0} % triangle left corner y=ct coordinate in S
  \pgfmathsetmacro\Lz{4*\d} % proper/rest length L0 in S
  \pgfmathsetmacro\L{\Lz/cos(\ang)} % length L in S'
  \coordinate (L) at (\Lz,0); % rod end in S
  \coordinate (L') at (\ang:\L); % rod end in S'
  \coordinate (A) at (\xA,\yA); % point A in triangle
  \coordinate (B) at (\xA+\Lz,\yA); % point B in triangle
  \coordinate (B') at ($(B)+(90-\ang:0.61)$); % point B' in triangle
  \coordinate (B'') at ($(B)+(180+\ang:\L+0.25)$); % point B' in triangle
  % FILL
  \pgfmathsetmacro\Lmax{1.8*\L}
  \begin{pgfonlayer}{back} % draw on back
    \fill[mylightblue2] (270-\ang:\Lmax) --++ (90-\ang:2*\Lmax) -- ++ (0:\Lz) --++ (270-\ang:2*\Lmax) -- cycle;
  \end{pgfonlayer}
  \draw[->,thick,mygreen] (0,0) --++ (90-\ang:\Lmax);
  \draw[thick,mygreen] (0,0) --++ (270-\ang:\Lmax);
  \draw[->,thick,myblue] (\Lz,0) --++ (90-\ang:\Lmax);
  \draw[thick,myblue] (\Lz,0) --++ (270-\ang:\Lmax);
  
  % ROD
  \pgfmathsetmacro\el{0.4*\Lz/(1-tan(\ang))}
  \draw[rod] (\xA,0) --++ (L)
    node[pos=0.5,below=-1] {$L$};
  \draw[rod] ({\xA-\el*tan(\ang)},{-\el}) --++ (L) coordinate (E1);
  % \draw[rod] (\ang:{\xA/cos(\ang)}) --++ (L')
    % node[pos=0.485,above=1] {\contour{mylightblue2}{$L'$}};
  
    %  LIGHT CONE
    \draw[photon] (0,0) -- (45:1.2*\Lmax);
    \draw[photon] (0,0) -- (135:1.2*\Lmax);
    \draw[photon] (0,0) -- (225:1.2*\Lmax);
    \draw[photon] (0,0) -- (315:1.2*\Lmax);
  % TRIANGLE
  % \draw[very thick,myred,rounded corners=0.1]
    % (A) -- (B') node[midway,above=0] {\contour{mylightblue2}{$\Delta x'$}}
        % -- (B) node[midway,right=-2] {\contour{myfieldred}{$c\Delta t$}}
        % -- cycle node[pos=0.51,below=-1] {\contour{mylightblue2}{$\Delta x$}};
  %\fill[myfieldred] (A)++(185:0.089) circle(0.04);
  %\fill[mydarkred] (A) circle(0.03) node[below=1,left=-2.7] {A};
  % \fill[myfieldred] (A)++(200:0.1) circle(0.04);
  % \fill[mydarkred] (A) circle(0.03) node[below=1,left=-2.3] {\contour{myfieldred}{A}};
  \fill[mygreen] (0,0) circle(0.03) node[left=5] {\contour{white}{$E_2$}};
  % \draw[mygreen,dashed,thick] (0,0) -- (0,-\ell);
  % \fill[mygreen] (0,-\ell) circle(0.03);
  % \draw[<->] (0,-\el) -- ++ (\el,0) node[midway,below] {$\ell$};
  \draw[<->] (0,-\el-0.05) -- ++ (\el,0) node[midway,below] {$\ell$};
  % \draw[<->] (0,-\el-0.05) -- ++ ({-\el*tan(\ang)},0) node[midway,below] {$\beta\ell$};
  % \fill[mydarkred] (B) circle(0.03) node[below=1,right=0] {\contour{myfieldred}{B}};
  % \draw[dashed,myblue] (\Lz,0) -- (B') node[above] {\contour{myfieldred}{$x_2'$}};
  % \draw[dashed,myblue] (\Lz,0) -- (B'') node[left] {\contour{myfieldred}{$ct_2'$}};
  \fill[myblue] (E1) circle(0.03) node[above=1,right=0] {\contour{white}{$E_1$}};
  \coordinate (E3) at ($(B)+(0,0.46)$);
  % \draw[-,mygreen] (90-\ang:-0.2*\Lmax) arc(-90-\ang:-90:0.2*\Lmax) node[pos=0.4,below=1] {\contour{white}{$\alpha$}};
  % \fill[myblue] (E3) circle(0.03) node[above left] {\contour{myfieldred}{$E_3$}};
  % \fill[mydarkred] (B') circle(0.03) node[above=2,right=-1] {\contour{myfieldred}{B$'$}};
  % \fill[mydarkred] (B') circle(0.03) node[above=2,right=-1] {\contour{myfieldred}{B$'$}};
  
  \coordinate (E3) at ($(B)+(0,0.46)$);
  \fill[myblue] (E1) circle(0.03) node[right=0] {\contour{white}{$E_2$}};
  \draw[myblue,dashed] (E1) -- (0,\el) node[left] {$ct_2$};
  \draw[myblue,dashed] (E1) -- (\el,0) node[below] {$x_2$};
  \draw[<->] (0,-\el-0.05) -- ++ (-{\el*tan(\ang)},0) node[midway,below] {\contour{white}{$vt_2$}};

\end{tikzpicture}




\end{document}